\documentclass[12pt,letterpaper]{amsart}
\usepackage[T1]{fontenc}
\IfFileExists{newtxtext.sty}{\usepackage{newtxtext,newtxmath}}{\usepackage{mathptmx}}
\usepackage[letterpaper,margin=1in]{geometry}
\usepackage{amsmath,amsthm}
\usepackage{microtype}
\usepackage[hidelinks]{hyperref}
\DeclareMathOperator{\vol}{vol}
\DeclareMathOperator{\Norm}{N}
\newtheorem{lemma}{Lemma}[section]
\newtheorem{proposition}[lemma]{Proposition}
\newtheorem{remark}[lemma]{Remark}
\title[Independent shell sums for $G_2$]{Independent valuation-shell calculations for the $G_2$ Shalika-germ revision}
\author{ChatGPT (OpenAI)}
\date{November 15, 2026\\Research note, not a manuscript revision}
\begin{document}
\begin{abstract}
This AI-aided note independently derives several valuation-shell expressions used in Draft~6 of the revision of D. Joyner's 1994 manuscript on transfer of Shalika germs for $G_2$. It addresses the arithmetic portion of Open Problem~(1)(c) and the shell-sum portion of Open Problem~(1)(b). The calculated meromorphic continuations agree with the stated factors $C'(q)$ and the proposed magnitude $|C_3(q)|$. The analysis identifies the isolated one-dimensional cubic factors with the local Langlands--Shelstad principal-value construction, but it does \emph{not} establish compatibility of the proposed regulator on the whole cubic cone, the geometric factorization of the cubic fiber integral, Lemma~5.14, or the exceptional-divisor proviso of Lemma~0.8. In particular, $(X_3)$ remains unproved by this note.
\end{abstract}
\maketitle

\section{Purpose and notation}
Let $F$ be a nonarchimedean local field with residual characteristic at least $5$, residue-field cardinality $q$, integer ring $\mathcal O$, and maximal ideal $\mathfrak p=(\varpi)$. Normalize additive Haar measure by $\vol(\mathcal O)=1$ and $|\varpi|=q^{-1}$. The shell $v_F(x)=k$ has measure $(1-q^{-1})q^{-k}$. We use the elementary meromorphic geometric series
\begin{equation}\label{eq:geom}
G(z)=\frac{z}{1-z};\qquad \sum_{k\geq1}z^k=G(z),
\end{equation}
initially only in its domain of convergence.

The valuation cells are evaluated separately, each initially in its own region of absolute convergence and then by meromorphic continuation. \emph{One global convergence region need not exist.} In particular, the one-parameter cubic integral $B_1(\sigma)$ below has no common absolute-convergence strip. Section~\ref{sec:pv} checks the isolated one-dimensional continuations against the local principal-value definition in \cite[\S1]{LS83}; the corresponding global cubic-cone calculation remains open. The labels refer to Theorem~7.6H and Claim~8.3H of the GitHub manuscript \cite{Jo26}. The attached independent script checks the rational identities exactly; it does not import the manuscript's checking scripts.

\section{Quadratic factors in Theorem 7.6H}
\subsection*{Ordinary convergent factor}
For a nonsquare $d\in F^\times$ set
\begin{equation}\label{eq:Is}
I_s(d)=\int_F\frac{ds}{|s^2-d|}.
\end{equation}
If $d$ is a nonsquare unit, then $|s^2-d|=1$ on $\mathcal O$ and equals $|s|^2$ outside it. Thus $I_s=1+q^{-1}=(q+1)/q$. If $v_F(d)=1$, the regions $s\in\mathcal O^\times$, $s\in\mathfrak p$, and $|s|>1$ contribute $1-q^{-1}$, $1$, and $q^{-1}$ respectively. Hence $I_s=2$. These evaluations involve no regularization.

\subsection*{A divergent one-dimensional tail}
For $\Re z>1$,
\begin{equation}\label{eq:tail}
\int_{|t|>1}|t|^{-z}\,dt
=(1-q^{-1})\sum_{k=1}^{\infty}q^{k(1-z)}
=(1-q^{-1})G(q^{1-z}).
\end{equation}
The value of its meromorphic continuation at $z=0$ is $-1$. This explains the number used for $\mathrm{PV}\int_{|t|>1}dt$ in the manuscript, without yet identifying this continuation with the specific global principal value.

\subsection*{The quadratic-character factor}
Put $L=F(\sqrt{d_\alpha d_\zeta})$, and let $(\cdot,\cdot)_F$ denote the quadratic Hilbert symbol. Under the continuation in \eqref{eq:tail} the expression is
\begin{equation}\label{eq:It}
I_t=\int_{\mathcal O}\bigl[(t^2-d_\alpha d_\zeta,d_\alpha)_F-1\bigr]\,dt.
\end{equation}
Choose representatives for the two quadratic square classes with valuations $0$ or $1$. If $L/F$ is ramified, the integrand is $0$ on $\mathcal O^\times$ and $-2$ on $\mathfrak p$. Thus $I_t=-2/q$. If $L/F$ is unramified in the biquadratic situation, both $d_\alpha$ and $d_\zeta$ have valuation $1$, and their product is a nonsquare unit times $\varpi^2$. For $t=\varpi y$ the character reduces to the Legendre symbol of $y^2-\varepsilon$, where $\bar\varepsilon$ is a nonsquare in the residue field. The identity
\[
\sum_{y\in\mathbb F_q}\left(\frac{y^2-\varepsilon}{\mathbb F_q}\right)=-1
\]
then gives $I_t=-(q+1)/q^2$. One proof of the identity counts solutions of $x^2-y^2=-\varepsilon$: there are $q-1$ solutions, since $(x-y)(x+y)=-\varepsilon$.

\subsection*{The meromorphic projective-line factor}
Let $K_\zeta=F(\sqrt d)$. In coordinates $\xi=x+y\sqrt d$, the form
\begin{equation}\label{eq:quadratic-form}
\Omega_2=\frac{dx\wedge dy}{y^2}
\end{equation}
is invariant, up to sign, under $\xi\mapsto\xi^{-1}$. This follows from $y'=-y/(x^2-dy^2)$ and the Jacobian $(x^2-dy^2)^{-2}$ in absolute value. Partition $\mathbb P^1(K_\zeta)$ into the affine pieces $\mathcal O_{K_\zeta}$ and $\xi^{-1}\in\mathfrak p_{K_\zeta}$. For $A,B\ge0$ define
\begin{equation}\label{eq:Zab}
Z_{A,B}(z)=\int_{x\in\mathfrak p^A\!,\ y\in\mathfrak p^B}|y|^{z-2}\,dx\,dy
=q^{-A}\frac{(1-q^{-1})q^{B(1-z)}}{1-q^{1-z}}.
\end{equation}
For an unramified quadratic extension, $\mathcal O_K=\mathcal O\oplus\mathcal O\sqrt d$ and $\mathfrak p_K=\mathfrak p\oplus\mathfrak p\sqrt d$. For a tamely ramified quadratic extension with $d=\varpi u$, the second lattice is $\mathfrak p\oplus\mathcal O\sqrt d$. Therefore the two continuations are
\begin{align}\label{eq:Ixi}
I_\xi^{\rm unr}&=\left.(Z_{0,0}+Z_{1,1})\right|_{z=0}=-\frac2q,\\
I_\xi^{\rm ram}&=\left.(Z_{0,0}+Z_{1,0})\right|_{z=0}=-\frac{q+1}{q^2}.
\end{align}
Combining \eqref{eq:Is}, \eqref{eq:It} and \eqref{eq:Ixi} with $|d_\alpha d_\zeta|^{-1/2}$ gives
\begin{equation}\label{eq:Cprime}
|d_\alpha d_\zeta|^{-1/2}|I_tI_sI_\xi|=
\begin{cases}
4(q+1)q^{-5/2},& L/F\text{ ramified},\\
2(q+1)^2q^{-3},&L/F\text{ unramified}.
\end{cases}
\end{equation}
The two ramified-$L$ distributions of the ramification between $d_\alpha$ and $d_\zeta$ give the same constant.

\section{The cyclic cubic projective-line factor}
Let $K/F$ be cyclic cubic. For $\eta\in K$ set $\eta_i=\tau^i\eta$ ($i=0,1,2$) and
\[
D(\eta)=(\eta_0-\eta_1)(\eta_1-\eta_2)(\eta_2-\eta_0),\qquad
\Omega_3=\frac{d\eta_0\wedge d\eta_1\wedge d\eta_2}{D(\eta)}.
\]
The form $\Omega_3$ is unchanged, up to sign, under $\eta\mapsto\eta^{-1}$: numerator and denominator both acquire a factor $\Norm_{K/F}(\eta)^{-2}$. As in Section~2, use the two disjoint charts $\eta\in\mathcal O_K$ and $\eta^{-1}\in\mathfrak p_K$.

\subsection*{Unramified cubic extension}
Choose an integral power basis $\eta=a+b\theta+c\theta^2$ whose reduction generates $\mathbb F_{q^3}$ over $\mathbb F_q$. The determinant of the three embeddings is a unit. Moreover $D(\eta)$ is a unit precisely when $(b,c)$ is primitive. For arbitrary $(b,c)$, $v_F(D)=3\min(v_F(b),v_F(c))$. Inserting a local Mellin weight of order $z$ along $D$, the first chart therefore gives
\[
Z_{\rm unr}(z)=\frac{1-q^{-2}}{1-q^{1-3z}}.
\]
The inverted chart requires $a,b,c\in\mathfrak p$; its contribution is
$q^{-1}(1-q^{-2})q^{1-3z}/(1-q^{1-3z})$. Evaluating at $z=0$, each contributes $-(q+1)/q^2$.

\subsection*{Tamely ramified cubic extension}
Now $q\equiv1\pmod3$. Choose $\theta^3=\varpi u$ with $u\in\mathcal O^\times$ and $\mathcal O_K=\mathcal O[\theta]$. A primitive cube root of unity is in $F$. For $\eta=a+b\theta+c\theta^2$, direct multiplication of the three conjugate differences gives, up to a unit,
\[
D(\eta)=\varpi u\,(b^3-\varpi u c^3).
\]
The numerator determinant has the same initial factor $\varpi u$ in absolute value. Hence the measure density is $|b^3-\varpi u c^3|^{-1}da\,db\,dc$. The terms of the primitive $(b,c)$-shell are as follows: $b$ a unit contributes measure $1-q^{-1}$ and denominator norm $1$; $b\in\mathfrak p$ and $c$ a unit contributes measure $q^{-1}(1-q^{-1})$ and denominator norm $q^{-1}$. Subsequent simultaneous dilations by $\varpi^k$ multiply the integral by $q^{k(1-3z)}$. Thus
\[
Z_{\rm ram}(z)=\frac{(1-q^{-1})(1+q^{-z})}{1-q^{1-3z}},\qquad
Z_{\rm ram}(0)=-\frac2q.
\]
Since $\mathfrak p_K=\mathfrak p\oplus\mathcal O\theta\oplus\mathcal O\theta^2$, the inverted chart contributes $q^{-1}Z_{\rm ram}(0)=-2/q^2$. We have independently derived the formula asserted in Claim~8.3H:
\begin{equation}\label{eq:Iproj}
I_{\mathbb P^1(K)}=-\frac{2(q+1)}{q^2}
\quad\text{for both unramified and tamely ramified cyclic cubic }K/F.
\end{equation}

\section{The two cubic character integrals}
Let $\chi=\omega_{K/F}$ be the order-$3$ character attached to the cyclic cubic extension. The proof sketch introduces
\begin{equation}\label{eq:Bdefs}
B_1(S)=\int_F\chi(y)|y|^{-\sigma}|y+1|^{4\sigma}\,dy,
\qquad S=q^\sigma,
\end{equation}
and
\[
B_2(1)=\operatorname{PV}\!\int_F\chi(3u-1)
\frac{|3u-1|}{|u-1|^3}\,du.
\]
The following are calculations of the \emph{specified} meromorphic shell sums. At $u=1$ the pole of $B_2$ is regularized by multiplying the local density by $|u-1|^z$ before continuation. These are local regularizations, whose compatibility with \cite[\S1]{LS83} is checked in Section~\ref{sec:pv}. Importantly, $B_1(\sigma)$ has \emph{no} global region of absolute convergence: convergence near $y=-1$ demands $\operatorname{Re}\sigma>-1/4$, whereas at infinity it requires $\operatorname{Re}\sigma<-1/3$.

\subsection*{Unramified cubic character}
Let $w=\chi(\varpi)$, so $w^2+w+1=0$. Split $B_1$ into $v(y)>0$, $v(y)=0$ with $\bar y\neq-1$, $y\in-1+\mathfrak p$, and $v(y)<0$, respectively. The four contributions give
\begin{equation}\label{eq:B1u}
B_1(S)=(1-q^{-1})\left[G\!\left(\frac{wS}{q}\right)
+G\!\left(\frac1{qS^4}\right)
+G(w^2qS^3)\right]+1-\frac2q.
\end{equation}
At $S=1$, use $w^3=1$ to see that the first and third terms cancel the other contributions. Thus $B_1(1)=0$, and direct differentiation gives
\begin{equation}\label{eq:B1up}
B_1'(1)=4(q-1)\left[\frac{w}{(q-w)^2}-\frac1{(q-1)^2}\right].
\end{equation}
For $B_2$, put $y=u-1$, so $3u-1=3y+2$. The shell $y\in\mathfrak p$ gives $-q/(q+1)$ by continuing $\int_{\mathfrak p}|y|^{z-3}dy$. Units outside $y\equiv-2/3\pmod{\mathfrak p}$ contribute $1-2/q$; the exceptional unit class contributes $(1-q^{-1})G(w/q^2)$; and $v(y)<0$ contributes $(1-q^{-1})G(w^2/q)$. Therefore
\begin{equation}\label{eq:B2u}
B_2(1)=-\frac q{q+1}+1-\frac2q
+(1-q^{-1})\left[G\!\left(\frac w{q^2}\right)+G\!\left(\frac{w^2}{q}\right)\right].
\end{equation}
Multiply \eqref{eq:B1up} by \eqref{eq:B2u} and reduce using $w^2+w+1=0$. The result is the exact identity
\begin{equation}\label{eq:unramproduct}
B_1'(1)B_2(1)=\frac{12}{q^2-1}.
\end{equation}

\subsection*{Tamely ramified cubic character}
Now $\chi$ is a nontrivial cubic character on the residue-field units. Shells with $v(y)\neq0$ in $B_1$ vanish by multiplicative character orthogonality. The class $y\equiv-1$ gives the geometric series, and the other unit classes sum to $-1/q$. Hence
\begin{equation}\label{eq:B1r}
B_1(S)=-\frac1q+(1-q^{-1})G\!\left(\frac1{qS^4}\right),
\quad B_1(1)=0,\quad B_1'(1)=-\frac4{q-1}.
\end{equation}
In $B_2$, shells with $v(y)<0$ and the exceptional class $y\equiv-2/3$ have zero character average. For the ordinary unit residues the finite character sum is $-\chi(2)/q$. On $y\in\mathfrak p$, the continued value is $-\chi(2)q/(q+1)$. Thus
\begin{equation}\label{eq:B2r}
B_2(1)=-\chi(2)\left(\frac q{q+1}+\frac1q\right)
=-\chi(2)\frac{q^2+q+1}{q(q+1)}.
\end{equation}

Under the additional \emph{unproved} geometric and principal-value identifications used in Claim~8.3H, the stated factor
$-\frac{q-1}{4q}B_1'(1)B_2(1)$ has absolute value
$3/[q(q+1)]$ in the unramified case and $(q^2+q+1)/[q^2(q+1)]$ in the ramified case. Multiplying by \eqref{eq:Iproj} then gives precisely the \emph{proposed} magnitudes $6/q^3$ and $2(q^2+q+1)/q^4$ respectively. This is not a proof that the geometric fiber integral equals that product.


\section{Langlands--Shelstad principal values for the one-dimensional factors}\label{sec:pv}
We now compare the \emph{isolated} factors $B_1$ and $B_2$ with the primary
source \cite[\S1, Lemmas 1.A and 1.C]{LS83}, which was consulted for this
revision of the note. The comparison is local on the \emph{projective line};
it does not settle the push-forward of the cubic-cone fiber measure.

\begin{proposition}\label{prop:local-pv}
Let $\chi:F^\times\to\mathbb C^\times$ be a nontrivial character of order $3$.
For $\sigma$ in a neighborhood of $0$, let
\[
 h_\sigma(y)=\chi(y)|y|^{-\sigma}|1+y|^{4\sigma}
 \quad(y\in\mathbb P^1(F)\setminus\{0,-1,\infty\}).
\]
The Langlands--Shelstad principal value
$\oint_{\mathbb P^1(F)}h_\sigma(y)|dy|$ is well-defined and holomorphic
near $\sigma=0$. It equals the rational shell continuation denoted
$B_1(q^\sigma)$ in \eqref{eq:B1u} and \eqref{eq:B1r}. In particular,
\[
 \oint_{\mathbb P^1(F)}\chi(y)|dy|=0,
 \qquad B_1(1)=0.
\]
Likewise, the principal value on $\mathbb P^1(F)$ of
\[
 \chi(3u-1)\frac{|3u-1|}{|u-1|^3}\,|du|
\]
exists in the sense of \cite[\S1]{LS83} and equals the cellwise values
$B_2(1)$ calculated in \eqref{eq:B2u} and \eqref{eq:B2r}.
\end{proposition}

\begin{proof}
Use the disjoint compact-open covering of $\mathbb P^1(F)$ by the sets
$y\in\mathfrak p$, $y\in-1+\mathfrak p$, the remaining units, and
$z=1/y\in\mathfrak p$ (including the point at infinity).
In these coordinates the local character and measure factors are,
respectively,
\[
\begin{array}{c|c|c}
 \text{neighborhood}&\text{character}&\text{power of the local coordinate in }|dy|\,h_\sigma\\
 \hline
 y=0 &\chi(y)&1-\sigma\\
 y=-1 &1&1+4\sigma\\
 z=0\ (y=\infty)&\chi^{-1}(z)&-1-3\sigma
\end{array}
\]
where powers are those in the standard expression
$|u|^c\,|du|/|u|$. Near ordinary units there is no singularity.
Since $\chi$ and $\chi^{-1}$ are nontrivial and $1+4\sigma\ne0$ near
$\sigma=0$, the nonresonance condition (1.4) of \cite{LS83} is satisfied.
Lemma 1.A there gives, on each compact-open neighborhood, exactly the
geometric-series continuation used in Section~4. Their sum is holomorphic
near $\sigma=0$ and is the stated $B_1(q^\sigma)$. Equivalently,
$B_1(1)=0$ follows immediately from \cite[Lemma 1.C]{LS83} applied to
$\theta_0=\chi$, $\theta_1=\chi^{-1}$, $c_0=1$, $c_1=-1$.

For $B_2$, set $y=u-1$. The singular points are $y=0$, $y=-2/3$ and
$y=\infty$. Their local powers and characters are
\[
\begin{array}{c|c|c}
 y=0&1&-2\\
 v=3y+2=0&\chi(v)&2\\
 z=1/y=0&\chi^{-1}(z)&1
\end{array}
\]
up to locally constant unit factors. None violates the nonresonance
condition of \cite[(1.4)]{LS83}; $3$ is a unit since the residual
characteristic is at least $5$. The four valuation regions used to derive
\eqref{eq:B2u} and \eqref{eq:B2r} therefore give precisely the local
principal-value integrals by \cite[Lemma 1.A]{LS83}. This proves the claim.
\end{proof}

\begin{remark}\label{rem:global-pv}
The result does \emph{not} identify the quantity
$\oint_{\mathcal C_K(F)}\kappa_1(t_*)|\omega_{\mathcal Q_3}|$
with the product of the isolated principal values. That would require
showing that the factorization, compactification and push-forward of the
measure meet the smooth-submersion and local-coordinate hypotheses of
\cite[Lemma 1.F]{LS83}, or proving an appropriate extension across the
singular and boundary loci. In particular, the proposed radial regulator
$R_{\mathrm{cone}}^{\sigma/3}$ in Claim 8.3H has not been identified with
the independent local regulators of the normal-crossings construction.
The order of continuation and fiber integration cannot simply be interchanged.
\end{remark}

\section{Status and remaining tasks}
The arguments above independently derive the shell sums; Proposition~\ref{prop:local-pv} further identifies the isolated one-dimensional cubic factors with Langlands--Shelstad local principal values. The symbolic companion script reproduces the rational identities without importing any $G_2$ verification modules. This does \emph{not} establish the analytic fiber-integral identity required in the main draft.

Three essential tasks remain. First, prove that the global cubic-cone regulator and fiberwise factorization implement the principal-value construction of \cite{LS83}; the one-dimensional result of Proposition~\ref{prop:local-pv} is not enough. Second, justify the cone factorization and changes of variables in the sketch of Claim~8.3H, including all Jacobians, $X$-dependent factors, and boundary contributions. Third, retain Lemma~5.14 and the divisor condition of Lemma~0.8 as hypotheses until independently resolved. The $r=1$ comparison and the omitted torus classes are also untouched here.

\section*{Glossary}
\begin{description}
\item[Valuation shell.] The measurable subset $\{x\in F:v_F(x)=k\}$.
\item[Mellin continuation.] Continuation in a complex exponent $z$ from a region in which a weighted local integral converges.
\item[Principal value.] A regularized integral defined by a specified analytic construction; it must not be identified with an arbitrary meromorphic finite part without proof.
\item[Tamely ramified.] A finite extension whose ramification degree is prime to the residue characteristic.
\end{description}

\section*{Acknowledgments}
This note was written and algebraically verified by ChatGPT (OpenAI) on November 15, 2026, while reviewing the current GitHub files. It is a separate computational and mathematical note, not a revision of the manuscript by D.~Joyner, Claude, and Gemini. AI review is not an independent certification of the results. Further details of the project-level AI usage appear in the repository revision log.

\begin{thebibliography}{99}
\bibitem[Jo26]{Jo26} D.~Joyner, Claude, and Gemini, \emph{On the transfer of Shalika germs for $G_2$}, Draft~6 (2026), \url{https://github.com/wdjoyner/trace-formulas/tree/main/g2stuff}.
\bibitem[LS83]{LS83} R.~Langlands and D.~Shelstad, ``On principal values on $p$-adic manifolds'', in \emph{Lie Group Representations II}, Lecture Notes in Math. \textbf{1041}, Springer, 1984, 250--279. Original PDF consulted directly for Section~5: \url{https://www.sunsite.ubc.ca/DigitalMathArchive/Langlands/pdf/prin2-ps.pdf}.
\end{thebibliography}
\end{document}